\section{引言：为什么需要扩散语言模型}

自回归（Autoregressive, AR）模型长期以来是语言生成领域的主流范式。从 GPT 系列到 LLaMA，这些模型以严格的从左到右（left-to-right）方式逐词生成文本，已经取得了巨大的成功。然而，AR 模型存在一个根本性的速度瓶颈：生成 $N$ 个 token 需要 $N$ 步前向传播，无法实现真正的并行生成。

\begin{importantbox}{扩散语言模型的核心优势}
扩散语言模型（Diffusion Language Models）摆脱了从左到右的约束，能够以任意顺序生成文本，最重要的是可以\textbf{并行生成多个词}。这使得它们在推理速度上具有显著的优势。
\end{importantbox}

2025 年以来，扩散语言模型引起了工业界的广泛关注。Google 发布了 Gemini Diffusion，ByteDance 推出了 Seed Diffusion，Inception Labs 发布了 Mercury Coder。这些模型不仅性能出色，而且在推理速度上远超传统 AR 模型。Forbes 和 Fortune 等主流商业媒体也开始报道这一技术变革。

本次客座讲座的主讲人 Subham Sahoo 是 Cornell 大学博士毕业生，其博士研究聚焦于扩散语言模型。他是 MDLM（Masked Diffusion Language Models）和 SoLM（Sorting-based Language Models）等重要工作的第一作者。本讲座将深入讲解这两篇论文的核心思想与技术细节。

\subsection{扩散语言模型发展时间线}

\begin{knowledgebox}{从 2022 到 2025：关键里程碑}
\begin{enumerate}
    \item \textbf{2022年}：扩散语言模型首次被提出。这些模型能以随机顺序生成词，但文本质量远落后于 AR 模型，社区反应冷淡。
    \item \textbf{2024年初}：Lou et al. 证明扩散 LM 搭配现代 Transformer 架构可以匹配 GPT-2 级别的 AR 模型质量，同时显著更快。
    \item \textbf{2024年中}：MDLM 发布，重新推导并简化了离散扩散训练算法，在小规模文本数据集上匹配 AR 模型性能。
    \item \textbf{2025年初}：Diffusion Duality 论文提出蒸馏算法，使模型能用仅 10 步生成 1000 词的序列，质量几乎无损。
    \item \textbf{2025年}：SoLM 发布，融合 AR 与离散扩散，支持 KV 缓存，推理速度实现数量级提升。
\end{enumerate}
\end{knowledgebox}

\subsection{本章小结}

扩散语言模型代表了文本生成范式的一次重要转变。从 2022 年的概念提出到 2025 年的工业化部署，这一领域经历了快速发展。MDLM 框架是当前多个工业级扩散语言模型（如 ByteDance 的 Seed Diffusion）的基础。

\newpage
